\chapter{Empirical Results \& Analysis}
\label{ch:results}

This chapter reports every experiment in the order the argument requires.
Section~\ref{sec:res_phase0} validates the learning implementations on
tasks whose correct answer is known. Sections~\ref{sec:res_calibration} to
\ref{sec:res_ladder} then use the validation year to settle three things
that must be settled before the held-out months are spent: how coarse the
action model may be, whether the learners use their observations at all,
and whether each rung of the state ladder changes what is learned.
Section~\ref{sec:res_final} reports the held-out test months, evaluated
once. Section~\ref{sec:res_analysis} draws the findings together.

All numbers are produced by committed scripts from cached data; none are
illustrative. Complete per-month results appear in
Appendix~\ref{app:permonth}.

For a reader keeping only one question in mind, let it be this one:
\emph{did any learned agent do better than simply buying at the start of
the month and holding to the end?} The answer is no, and the interesting
part of this chapter is the evidence for \emph{why} not, which turns out to
be a property of the problem as posed rather than a failure of the
implementations.

\section{Experimental setup}
\label{sec:res_setup}

\paragraph{Data and split.} SPY daily auto-adjusted OHLCV bars,
January 2018 to December 2025, one calendar month of daily bars per
episode, 96 episodes split chronologically into 60 training months
(2018--2022), 12 validation months (2023) and 24 held-out test months
(2024--2025). Market features are computed on the continuous series before
the episodes are cut (Section~\ref{sec:data}).

\paragraph{Protocol.} Initial cash $C_0 = \$10{,}000$; proportional fee
$c = 5$ basis points; trade slice $\Delta = 0.25\,C_0$ as calibrated in
Section~\ref{sec:res_calibration}; 80{,}000 environment steps per method
per seed; seeds $\{42,\dots,47\}$; greedy masked evaluation. Metrics and
reference strategies are as defined in Section~\ref{sec:protocol}.

\paragraph{What the validation year is spent on.} Three decisions are taken
on 2023 and on nothing else: the trade slice, the reward scale, and where
the state ladder stops. Each is reported below with the rule that decided
it. The test months are read once, in Section~\ref{sec:res_final}.

\section{Phase 0: algorithm validation on games}
\label{sec:res_phase0}

Before any market experiment, the learning implementations were validated
on tasks with unambiguous success criteria
(Table~\ref{tab:res_phase0}). Every learner clears its random floor. Two
observations carry forward. First, from-scratch REINFORCE learns but
plateaus below the more stabilised methods on CartPole (292 vs.\ 500) ---
the textbook variance limitation of Monte-Carlo policy gradients,
visible in our own runs. Second, tabular Q-learning solves CartPole
\emph{only after} hand-discretizing the 4-dimensional continuous state;
the 9-dimensional continuous trading state makes any such table
infeasible, which is precisely why the value-based track uses a deep
Q-network (Section~\ref{sec:dqn}).

\begin{table}[h]
\centering
\caption{Phase 0 mean evaluation return (30 greedy episodes, seed 42).
CartPole solved $\approx 475$; LunarLander solved $\approx 200$.}
\label{tab:res_phase0}
\begin{tabular}{lrrr}
\toprule
\textbf{Algorithm} & \textbf{CartPole-v1} & \textbf{FlappyBird-v0} & \textbf{LunarLander-v3} \\
\midrule
Random                  & 28.8           & $-7.4$        & $-183.4$ \\
Tabular Q-learning      & \textbf{500.0} & ---           & --- \\
DQN                     & \textbf{500.0} & 7.0           & \textbf{177.6} \\
REINFORCE (scratch)     & 292.2          & 7.1           & 15.3 \\
A2C                     & \textbf{500.0} & 4.6           & $-41.2$ \\
PPO                     & \textbf{500.0} & \textbf{12.6} & 176.4 \\
\bottomrule
\end{tabular}
\end{table}

\begin{figure}[h]
  \centering
  \includegraphics[width=0.48\textwidth]{phase0_cartpole_comparison}
  \includegraphics[width=0.48\textwidth]{phase0_cartpole_curves}
  \caption{CartPole-v1: final comparison (left) and learning curves
  (right, 25-episode moving average). Scratch REINFORCE learns but
  plateaus below the stabilised methods --- the variance limitation
  predicted by the theory of Section~\ref{sec:pg_derivation}.}
  \label{fig:res_phase0_cartpole}
\end{figure}

\begin{figure}[h]
  \centering
  \includegraphics[width=0.48\textwidth]{phase0_flappy_comparison}
  \includegraphics[width=0.48\textwidth]{phase0_lunarlander_comparison}
  \caption{Flappy Bird (left) and LunarLander-v3 (right): every learner
  clears its random floor under a fixed budget.}
  \label{fig:res_phase0_others}
\end{figure}

The conclusion carried forward is narrow but necessary: the update rules
work. Any failure to learn on market data is therefore a statement about
the market problem, not about the code, and that distinction is what makes
the negative results below interpretable.

\section{Calibrating the action model}
\label{sec:res_calibration}

The first validation experiment involves no learning at all, and it
changed the experimental design.

Because a Buy commits a fixed fraction of initial capital, becoming fully
invested takes $C_0/\Delta$ purchases. A monthly episode is about 21 bars
long. The slice therefore does not merely set how finely the agent can
adjust its position: it caps the average exposure attainable within an
episode, and so caps the return attainable, regardless of skill. That
ceiling can be measured directly, by running the policy that accumulates as
fast as the action model allows, which bounds every policy at that
granularity.

\begin{table}[h]
\centering
\caption{Attainable exposure ceiling by trade slice, on the 12 validation
months (mean episode length 20.8 bars). The ceiling is measured by the
fastest legal accumulation policy, so no policy at that granularity can
exceed it. Buy-and-hold attains exposure 0.910 and $+\$166.92$.}
\label{tab:res_calibration}
\begin{tabular}{rrrrc}
\toprule
\textbf{Slice $\Delta$} & \textbf{Buys to full} &
\textbf{Ceiling exposure} & \textbf{Ceiling $\Delta W$} &
\textbf{Reaches target} \\
\midrule
$0.05\,C_0$ & 20 & 0.450 & $+\$95.49$  & no \\
$0.10\,C_0$ & 10 & 0.690 & $+\$121.79$ & no \\
$0.25\,C_0$ & 4  & 0.838 & $+\$167.39$ & \textbf{yes} \\
$0.50\,C_0$ & 2  & 0.886 & $+\$162.77$ & yes \\
\bottomrule
\end{tabular}
\end{table}

The consequence is stark. At $\Delta = 0.1\,C_0$, the slice used earlier in
this project, the best attainable mean exposure is 0.690 against
buy-and-hold's 0.910, and the best attainable monthly profit is
$+\$121.79$ against buy-and-hold's $+\$166.92$. An agent at that
granularity cannot beat buy-and-hold in a rising market however well it
trades. Any earlier claim that an agent ``failed to beat the benchmark''
under that setting was therefore reporting arithmetic as though it were a
result about learning.

Applying the rule declared in Section~\ref{sec:actions} --- the smallest
slice whose exposure ceiling reaches 90\% of buy-and-hold's, here 0.819 ---
selects $\Delta = 0.25\,C_0 = \$2{,}500$. The criterion is stated on
exposure rather than on return because the exposure ceiling rises
monotonically with the slice while the return ceiling does not:
$0.50\,C_0$ holds more on average than $0.25\,C_0$ yet earns slightly less,
because a coarser slice buys at a different set of prices. Selecting on the
non-monotone quantity would make the choice depend on the particular path
2023 happened to take. The coarser slice is retained as a comparison row in
Section~\ref{sec:res_final}, so its cost is measured rather than merely
avoided.

\section{Do the learners use their state?}
\label{sec:res_conditioning}

The second validation experiment asks a question that returns cannot
answer. A policy that buys whenever a Buy is legal requires no market
information whatsoever, and in a rising market it performs well. Reward
alone cannot distinguish such a policy from one that has learned something,
so the conditioning diagnostic of Section~\ref{sec:protocol} is reported
alongside it.

The question was prompted by an anomaly. In the first ladder run, trained
agents at rungs from one feature to twelve returned monthly profits that
were not merely similar but identical to several decimal places, and
identical to a reference strategy's. Policies reading different
observations do not agree to that precision.

\begin{table}[h]
\centering
\caption{Conditioning diagnostic on the 12 validation months at the
nine-feature state, six seeds. ``State-dependent'' counts the seeds whose
policy chose different actions in states offering the same legal choices.
The reward scale is learner-side preprocessing only; the environment emits
dollars in both columns.}
\label{tab:res_conditioning}
\begin{tabular}{llrrc}
\toprule
\textbf{Method} & \textbf{Reward scale} & \textbf{Mean $\Delta W$} &
\textbf{Seed spread} & \textbf{State-dependent} \\
\midrule
REINFORCE & dollars              & $+\$100.64$ & \$121.71 & \textbf{0 / 6} \\
REINFORCE & fraction of capital  & $+\$106.24$ & \$85.51  & 1 / 6 \\
Deep Q    & dollars              & $+\$69.92$  & \$83.88  & \textbf{6 / 6} \\
Deep Q    & fraction of capital  & $+\$48.19$  & \$124.04 & 5 / 6 \\
\bottomrule
\end{tabular}
\end{table}

\paragraph{REINFORCE does not read its state.} In none of the six seeds
does the trained policy distinguish two states that offer the same legal
actions. It converges instead to one of two rules: buy whenever a Buy is
legal, or never trade. Its per-seed profits are
$\{+0.1,\, +119.2,\, +119.2,\, +121.8,\, +121.8,\, +121.8\}$, and
$+\$121.79$ is exactly the slice-schedule benchmark's figure. The nine
features are present, computed correctly and passed to the network, and the
network ignores them.

This is a consequence of the algorithm meeting this particular reward
rather than a defect. In a month whose price rises, the return following a
Buy exceeds the episode average more often than not, so the gradient step
raises the log-probability of Buy in whatever state it was taken. Repeated
across many episodes, this pushes the Buy logit up regardless of the state
that accompanied it, and the softmax saturates on a state-independent
action. Plain REINFORCE carries no baseline and no value function to
separate ``this action was good'' from ``this month was good''; the
variance limitation visible on CartPole in Section~\ref{sec:res_phase0}
appears here as a collapse to a constant policy.

The consequence for this dissertation is specific and it must be stated
plainly: \emph{REINFORCE cannot inform a comparison between state
representations}, because it does not use the representation. Its rows in
the next section are reported for completeness and interpreted only as
evidence about the algorithm.

\paragraph{Deep Q-learning does read its state.} In all six seeds the
learned policy chooses differently in states offering identical legal
actions. Its exposure varies within and across months rather than saturating
at a corner. It is therefore the method through which the ladder can be
read, and the fact that its mean profit is \emph{lower} than REINFORCE's is
the beginning of the substantive finding rather than a mark against it: it
is deviating from full exposure on the basis of what it observes, and in a
year the market spent rising, deviation costs money.

\paragraph{Reward scaling was tested and rejected.} A competing explanation
for the large seed spread was numerical conditioning: Q-values estimate a
sum of dollar rewards across a month, reaching the low hundreds, while the
network is initialised near zero, so a squared-error loss produces very
large early gradients. Dividing the reward by initial capital rescales the
targets to the order of $0.02$ without altering the MDP. It does not help.
REINFORCE is unaffected, as its within-episode standardisation of returns
predicts, and Deep Q-learning is made worse on all three measures: mean
profit falls from $+\$69.92$ to $+\$48.19$, the seed spread rises from
\$83.88 to \$124.04, and one seed loses state-dependence entirely. The
reward therefore remains in dollars exactly as
Equation~(\ref{eq:reward}) defines it, and the observed variance is a
property of the problem rather than of the units.

\section{The ladder: does more information help?}
\label{sec:res_ladder}

Chapter~3 argued the state ladder from information requirements, each rung
resolving a situation its predecessor could not distinguish. This section
tests whether those additions change what is learned. The reward and the
action model are held fixed throughout, so any difference between adjacent
rows is attributable to the information added between them.

\begin{table}[h]
\centering
\caption{The ladder on the 12 validation months, six seeds per cell, at the
calibrated slice. Spread is the range across seeds. ``St.-dep.'' counts
state-dependent seeds. Reference: B1a slice schedule $+\$167.39$,
buy-and-hold $+\$166.92$, random legal $+\$64.99$, do-nothing \$0.}
\label{tab:res_ladder}
\small
\begin{tabular}{llrrcrrcr}
\toprule
& & \multicolumn{3}{c}{\textbf{REINFORCE}} & \multicolumn{4}{c}{\textbf{Deep Q-learning}} \\
\cmidrule(lr){3-5}\cmidrule(lr){6-9}
\textbf{Rung} & \textbf{Dim.} & \textbf{Mean $\Delta W$} & \textbf{Spread} &
\textbf{St.-dep.} & \textbf{Mean $\Delta W$} & \textbf{Spread} &
\textbf{St.-dep.} & \textbf{Expo.} \\
\midrule
$S_0$   & 1  & $+161.18$ & 37.3  & 0/6 & $+68.68$  & 159.4 & 5/6 & 0.50 \\
$S_1$   & 3  & $+119.29$ & 125.7 & 0/6 & $-2.43$   & 16.4  & 5/6 & 0.06 \\
$S_2$   & 4  & $+161.18$ & 37.3  & 0/6 & $+64.47$  & 139.6 & 6/6 & 0.36 \\
$S_{3a}$& 5  & $+127.98$ & 124.6 & 1/6 & $+102.59$ & 28.8  & 6/6 & 0.56 \\
$S_3$   & 9  & $+154.96$ & 37.3  & 0/6 & $+86.16$  & 64.2  & 6/6 & 0.33 \\
$S_4$   & 10 & $+148.75$ & 37.3  & 0/6 & $+83.86$  & 155.2 & 6/6 & 0.38 \\
$S_5$   & 12 & $+163.88$ & 19.1  & 2/6 & $+40.20$  & 69.7  & 6/6 & 0.29 \\
\bottomrule
\end{tabular}
\end{table}

\paragraph{The declared stopping rule terminates at the first rung.} The
rule fixed in Section~\ref{sec:protocol} was to stop at the first rung
whose successor moves mean profit, intra-episode drawdown and trade count
by less than the across-seed spread. Applied to each learner separately and
to both pooled, it stops at $S_0$. No pair of adjacent rungs differs by
more than the variation between seeds of the same rung.

It is important to be precise about what that does and does not mean. It
does not mean a one-feature state is informationally sufficient; the
arguments of Section~\ref{sec:state_admissible} are unaffected by any
experiment, and $S_0$ remains a state in which the agent cannot tell
whether it holds a position. What it means is that the experiment lacks the
power to distinguish the rungs, because the outcome varies more with the
random seed than with the representation.

\paragraph{Why the variance is so large: bimodality, not noise.} The spread
is not the ordinary scatter of a noisy estimator. REINFORCE's per-seed
outcomes concentrate on a handful of values that recur across unrelated
rungs, because the policies it finds are corners: fully committed, or
inert. A mean between two corners describes neither, which is why the
seed count was raised from three to six and why the table reports the
spread rather than a standard error.

\paragraph{No rung beats an unconditional benchmark.} Every Deep Q row
falls below the slice schedule's $+\$167.39$, and its mean exposure sits
between 0.06 and 0.56 against the attainable 0.838. The agent is not
failing to trade; it is choosing to hold less than it could, on the basis of
features it demonstrably reads, in a year when holding more would have paid.
Adding information does not reverse this, and $S_5$, the richest state
tested, produces the lowest Deep Q profit of any rung.

\paragraph{Why this is the predicted outcome.}
Section~\ref{sec:reward} stated the reason before any of these runs: under
a risk-neutral objective on an asset with positive drift, the optimal policy
is unconditional maximum exposure. If the best available policy does not
depend on the state, then no state can be better than the smallest one, and
additional features can only add parameters to estimate and variance to the
estimate. The ladder does not fail to show an effect; it shows that the
\emph{objective} leaves no effect to find. An agent that discovers
buy-and-hold under this reward has solved the problem as posed, and the
richer states are solutions to a problem nobody asked it.

This is the sense in which the ladder is a design argument that survives its
own null result. The features are justified by what the agent must be able
to distinguish, and the experiment shows that distinguishing is not what
this reward pays for. It follows that the productive axis is the reward, not
the state, which is what the next section tests.

\section{Held-out results}
\label{sec:res_final}

The 24 test months, January 2024 to December 2025, are evaluated once. Every
configuration below was fixed before these months were read: the state ladder
stops where Section~\ref{sec:res_ladder} left it, the action model is the
calibrated slice, and the penalty weight was chosen on validation as described
in Section~\ref{sec:res_risk}.

\begin{table}[h]
\centering
\caption{Held-out test results, 24 months, six seeds per learned row. Spread is
the range across seeds. ``St.-dep.'' counts seeds whose policy reads its
observation. Drawdown is the mean intra-episode maximum. The three upper blocks
vary state and reward at the calibrated slice; the lower block varies the action
model at the frozen state and reward.}
\label{tab:res_final}
\small
\begin{tabular}{llrrrrrrc}
\toprule
\textbf{Configuration} & \textbf{Method} & \textbf{Mean $\Delta W$} &
\textbf{Spread} & \textbf{Expo.} & \textbf{Trades} & \textbf{Draw.} &
\textbf{Sharpe} & \textbf{St.-dep.} \\
\midrule
\multicolumn{9}{l}{\emph{Reference strategies (calibrated slice)}} \\
& B0 do-nothing      & $0.00$    & ---    & 0.000 & 0.0  & 0.0000 & 0.000 & --- \\
& B1a slice schedule & $+179.75$ & ---    & 0.841 & 5.0  & 0.0274 & \textbf{0.684} & --- \\
& B1b buy-and-hold   & $\mathbf{+180.25}$ & --- & 0.913 & 2.0 & 0.0310 & 0.638 & --- \\
& B2 random legal    & $+22.73$  & ---    & 0.304 & 13.3 & 0.0133 & 0.124 & --- \\
\midrule
\multicolumn{9}{l}{\emph{$S_3$, wealth-change reward} --- the baseline cell} \\
& REINFORCE & $+171.51$ & 24.7  & 0.810 & 10.3 & 0.0266 & 0.675 & 0/6 \\
& Deep Q    & $+74.68$  & 71.4  & 0.312 & 8.2  & 0.0124 & 0.507 & \textbf{6/6} \\
& PPO       & $+127.35$ & 179.8 & 0.598 & 3.7  & 0.0195 & 0.562 & 0/6 \\
\midrule
\multicolumn{9}{l}{\emph{$S_4$, wealth-change reward} --- adds drawdown to the state only} \\
& REINFORCE & $+167.39$ & 24.7  & 0.794 & 13.0 & 0.0263 & 0.670 & 0/6 \\
& Deep Q    & $+77.81$  & 89.5  & 0.354 & 7.4  & 0.0140 & 0.505 & \textbf{6/6} \\
& PPO       & $+159.97$ & 94.0  & 0.759 & 7.3  & 0.0250 & 0.671 & 1/6 \\
\midrule
\multicolumn{9}{l}{\emph{$S_4$, risk-aware reward}, $\lambda = 0.10$ --- adds drawdown to state and reward} \\
& REINFORCE & $+167.39$ & 24.7  & 0.794 & 13.0 & 0.0263 & 0.670 & 0/6 \\
& Deep Q    & $+52.90$  & 43.1  & 0.178 & 5.4  & \textbf{0.0082} & 0.416 & \textbf{6/6} \\
& PPO       & $+21.48$  & 128.9 & 0.108 & 0.7  & 0.0037 & 0.107 & 1/6 \\
\midrule
\multicolumn{9}{l}{\emph{Action model}, $S_3$ and wealth-change reward held fixed} \\
whole share & REINFORCE & $+77.91$  & 99.8  & 0.422 & 19.0 & 0.0151 & 0.427 & 0/6 \\
(B1a $+94.69$) & Deep Q & $+58.23$  & 61.6  & 0.177 & 10.4 & 0.0061 & 0.541 & \textbf{6/6} \\
            & PPO       & $+25.75$  & 90.1  & 0.128 & 4.8  & 0.0044 & 0.211 & 2/6 \\
\addlinespace
slice $0.1\,C_0$ & REINFORCE & $+117.40$ & 132.4 & 0.578 & 15.9 & 0.0196 & 0.550 & 0/6 \\
(B1a $+141.31$) & Deep Q & $+77.37$  & 42.9  & 0.327 & 11.2 & 0.0117 & 0.504 & \textbf{6/6} \\
            & PPO       & $+129.74$ & 69.4  & 0.635 & 10.0 & 0.0214 & 0.627 & 1/6 \\
\bottomrule
\end{tabular}
\end{table}

\paragraph{No learned agent beats buy-and-hold, on either measure.} The best
learned return is REINFORCE's $+\$171.51$ against buy-and-hold's $+\$180.25$,
and the best learned Sharpe ratio is $0.675$ against the slice schedule's
$0.684$. Neither gap is large, and neither has the wrong sign. Every method wins
in roughly a third of individual months, which is what a lower-exposure strategy
does in a market that falls in about a third of its months; it is not evidence
of timing.

\paragraph{The policy-gradient methods do not read their state.} REINFORCE is
mask-determined in all six seeds and PPO in all six at $S_3$, reproducing on the
test months exactly what Section~\ref{sec:res_conditioning} found on validation.
REINFORCE's near-benchmark return is therefore not a near-miss at beating the
market; it is the arithmetic of a policy that buys whenever a Buy is legal, which
is what B1a already is. Its apparent respectability and PPO's enormous seed
spread of \$179.75 have the same cause: both are choosing between corners rather
than conditioning on anything.

That the same collapse afflicts a modern stabilised method as well as the
textbook one is worth stating, because it rules out the most obvious
explanation. REINFORCE's variance alone cannot be the cause, since PPO's
clipped surrogate and value baseline exist precisely to control that variance
and it collapses anyway. What the two share is that they optimise a policy
directly against a reward whose best unconditional answer is already very good.

\paragraph{Deep Q-learning reads its state, and uses it to hold less.} It is
state-dependent in every seed of every cell, and it is also the weakest earner,
at $+\$74.68$ against $+\$180.25$. The two facts are one fact. Its mean exposure
is $0.312$ against the attainable $0.841$, and its drawdown is correspondingly
lower, $0.0124$ against buy-and-hold's $0.0310$.

The question that matters is whether that is risk management or merely
abstention, and it has a quantitative answer. Deep Q earns $41.4\%$ of
buy-and-hold's return while incurring $39.9\%$ of its drawdown. The ratio of
those two figures is $0.96$: the drawdown reduction exceeds the return sacrifice
by four percent, which is inside the seed spread on both quantities. Reducing
drawdown by holding less is something any investor can do without a neural
network, by holding less. The agent has not found a way to hold the same
exposure through smaller losses, and its Sharpe ratio, $0.507$ against
buy-and-hold's $0.638$, records that the trade was not worth making.

\paragraph{Adding drawdown to the state alone changes nothing.} Comparing the
first two blocks answers the question Chapter~3 deliberately left open: does
representing risk history help even when the reward does not ask for it? For
Deep Q the answer is no, $+\$74.68$ against $+\$77.81$ with a seed spread of
\$71 and \$89, so the difference is a small fraction of the noise. Exposure,
trade count, drawdown and Sharpe are likewise unchanged. PPO's figure rises from
$+\$127.35$ to $+\$159.97$, but with spreads of \$180 and \$94 that movement
carries no weight, and it is in any case a method that ignores the feature in
five of six seeds.

This is the cleanest vindication of the decision to keep drawdown out of the
core state. Had it been included by default, as originally proposed, this table
would have contained no comparison capable of showing that it earns nothing on
its own, and a reader would have had to take the inclusion on trust.

\paragraph{The action model matters more than the state.} The lower block is the
largest effect in the table. Holding state and reward fixed and changing only how
a Buy is sized moves REINFORCE from $+\$77.91$ to $+\$171.51$, a range wider than
anything the state axis produces. The explanation is the ceiling of
Section~\ref{sec:res_calibration}: under whole-share trading the fastest legal
accumulation reaches exposure $0.502$ and $+\$94.69$, and under the tenth-slice
model $0.694$ and $+\$141.31$. REINFORCE tracks those ceilings closely, at
$0.422$ and $0.578$ exposure, because tracking the ceiling is what a
buy-when-legal policy does.

Deep Q is the exception, and instructively so: it earns $+\$58.23$, $+\$77.37$
and $+\$74.68$ across the three action models, essentially flat. An agent that
has already decided to hold about a third of capital is indifferent to how
finely it could have held more. The action model constrains the policies that
saturate exposure and is irrelevant to the one that does not.

\section{The risk-aware reward}
\label{sec:res_risk}

\paragraph{Choosing the penalty weight, and a rule that had to be repaired.}
The weight $\lambda$ was selected on validation, and the first attempt to select
it failed in a way worth reporting rather than quietly correcting. The rule as
originally declared was to take the largest $\lambda$ whose mean wealth change
stayed within one across-seed spread of the unpenalised agent. On the validation
months that agent earned $+\$78.35$ with an across-seed spread of \$106.94, so
the tolerance exceeded the entire quantity it was meant to protect and the floor
fell to $-\$28.59$. Every candidate cleared it, including those that stopped
trading altogether, and the rule selected $\lambda = 2.0$ and an agent that never
opened a position in any month. An additive tolerance expressed as a multiple of
the noise is vacuous once the noise is larger than the signal.

The rule was therefore restated multiplicatively and with a participation
requirement: take the largest $\lambda$ retaining at least half the unpenalised
agent's mean wealth change \emph{and} at least half its mean exposure. The second
condition is what excludes the degenerate answer, because an agent holding
nothing has no drawdown, and reporting abstention as risk management would be a
category error. Both the superseded rule and its outcome are recorded in the
results file.

\begin{table}[h]
\centering
\caption{Penalty-weight sweep on the 12 validation months, Deep Q-learning, six
seeds. The unpenalised agent earns $+\$78.35$ at exposure $0.445$ with mean
intra-episode drawdown $0.0163$. Floors: $\Delta W \ge +\$39.17$ and exposure
$\ge 0.223$.}
\label{tab:res_lambda}
\begin{tabular}{rrrrrc}
\toprule
$\lambda$ & \textbf{Mean $\Delta W$} & \textbf{Exposure} &
\textbf{Drawdown} & \textbf{$\Delta$ drawdown} & \textbf{Clears floors} \\
\midrule
0.05 & $+88.78$ & 0.401 & 0.0128 & $-21.5\%$  & yes \\
0.10 & $+67.30$ & 0.344 & 0.0124 & $-24.0\%$  & \textbf{yes (selected)} \\
0.25 & $+19.90$ & 0.074 & 0.0035 & $-78.4\%$  & no \\
0.50 & $+1.73$  & 0.002 & 0.0001 & $-99.4\%$  & no \\
1.00 & $0.00$   & 0.000 & 0.0000 & $-100\%$   & no \\
2.00 & $0.00$   & 0.000 & 0.0000 & $-100\%$   & no \\
\bottomrule
\end{tabular}
\end{table}

The sweep shows why the participation condition is needed. The trade-off is
usable over a narrow range and then collapses: between $\lambda = 0.10$ and
$\lambda = 0.25$ exposure falls from $0.344$ to $0.074$ and the agent
substantially stops trading, and by $\lambda = 0.50$ it has stopped. The penalty
is $\lambda\,C_0$ per unit of fractional drawdown, so at $\lambda = 0.5$ a one
percent deepening of drawdown costs \$50 against a typical daily wealth change of
a few tens of dollars. Once the penalty outweighs the return term, the optimal
policy is to hold cash, and the agent finds it.

\paragraph{On the test months, the penalty buys nothing risk-adjusted.} At the
selected weight, Deep Q's drawdown falls from $0.0124$ to $0.0082$, a reduction
of $34\%$, and its return falls from $+\$74.68$ to $+\$52.90$, a reduction of
$29\%$. Its Sharpe ratio falls from $0.507$ to $0.416$. Measured against
buy-and-hold, the ratio of drawdown retained to return retained is $0.90$,
against $0.96$ for the unpenalised agent: a marginal improvement, well inside the
seed spread.

The honest reading is that the reward change works exactly as its algebra says it
must, and that this is not the same as working as intended. Penalising increases
in drawdown does reduce drawdown. It reduces it by reducing participation, which
is the cheapest way for the agent to comply, and it does not induce the behaviour
the experiment was looking for, which was the same exposure carried through
smaller losses. PPO's collapse to $0.7$ trades per month and $0.108$ exposure is
the same effect in more extreme form.

This is a negative result about a specific reward, not about risk-sensitive
reinforcement learning. A penalty on downside deviation gives the agent no credit
for return earned per unit of risk, so abstention always dominates. Objectives
that reward the ratio rather than penalising the numerator, such as the
differential Sharpe ratio of Moody and Saffell \cite{moody2001}, do not have this
degenerate solution, and Chapter~6 identifies that as the natural continuation.

\section{Fee sensitivity}
\label{sec:res_fees}

The 5-basis-point assumption is defensible for a retail investor in a highly
liquid instrument, but it is favourable, and a result that survives only at an
optimistic cost is not a result. The fee is therefore swept from frictionless to
deliberately punitive, on validation.

\begin{table}[h]
\centering
\caption{Fee sensitivity on the 12 validation months at the nine-feature state,
six seeds. Buy-and-hold trades twice per month, so its cost is nearly
fee-independent in relative terms; the agents trade far more.}
\label{tab:res_fees}
\begin{tabular}{rrrrrrr}
\toprule
\textbf{Fee} & \textbf{B1b $\Delta W$} & \textbf{REINF.\ $\Delta W$} &
\textbf{Trades} & \textbf{Deep Q $\Delta W$} & \textbf{Trades} &
\textbf{Exposure} \\
\midrule
0 bps  & $+177.09$ & $+177.56$ & 5.0  & $+108.00$ & 10.9 & 0.448 \\
5 bps  & $+166.92$ & $+154.96$ & 10.3 & $+86.16$  & 8.5  & 0.325 \\
10 bps & $+156.76$ & $+143.74$ & 4.7  & $+77.60$  & 8.9  & 0.436 \\
25 bps & $+126.33$ & $+88.76$  & 10.3 & $+32.35$  & 3.9  & 0.306 \\
50 bps & $+75.83$  & $+41.31$  & 2.8  & $+16.32$  & 1.9  & 0.181 \\
\bottomrule
\end{tabular}
\end{table}

\paragraph{Profit is fee-sensitive, as it must be.} Every strategy earns less as
costs rise, and the learned agents lose ground faster than buy-and-hold does,
because they trade between five and ten times as often. At 50 basis points Deep
Q retains $15\%$ of its frictionless profit while buy-and-hold retains $43\%$.
Nothing about that is surprising; it is arithmetic.

\paragraph{Behaviour is the informative measure, and it separates the methods.}
An agent that has represented the cost of trading should trade less as trading
becomes more expensive. Deep Q's trade count falls monotonically in the trend,
from $10.9$ per month at zero cost to $1.9$ at fifty basis points, a change of
$9.0$ against a maximum across-seed spread of $8.5$, and its exposure falls with
it. It has learned that trading is expensive.

REINFORCE has not. Its trade count over the same tenfold increase in cost reads
$5.0$, $10.3$, $4.7$, $10.3$, $2.8$: no trend, and the movement is smaller than
the seed noise of $15.8$. The pattern is exactly what the conditioning
diagnostic predicts, because the two corner policies it alternates between
involve about $5$ and about $21$ trades per month, so its average bounces
according to which corner each seed happened to reach. A policy that does not
read its state cannot read the cost of acting either.

This is the fee study's real contribution. Fees were placed inside the reward so
that the agent would pay them while learning
(Section~\ref{sec:reward}), and this table shows that placement working for one
method and having no behavioural effect on another, for a reason now
independently established.

\section{Analysis and key findings}
\label{sec:res_analysis}

\paragraph{1. The reward is the binding constraint, not the state.} Under
$r_t = \Delta W_t$ on an asset that drifts upward, the best policy is
unconditional maximum exposure, so no state can be better than the smallest and
extra features add only parameters and variance. This was stated in
Section~\ref{sec:reward} before any run, and three independent results now
support it: the ladder does not separate rungs from one feature to twelve
(Section~\ref{sec:res_ladder}); adding drawdown to the state without adding it to
the reward changes nothing beyond seed noise (Section~\ref{sec:res_final}); and
both direct policy optimisers converge on policies that ignore the state
entirely. The state design of Chapter~3 is not thereby refuted, because
admissibility is not an empirical question. What is refuted is the assumption
that a better representation is what this agent was missing.

\paragraph{2. Algorithm family determines whether the state is used at all.}
The two policy-gradient methods are mask-determined in almost every seed; the
value-based method is state-dependent in every seed of every cell. This is the
most robust behavioural finding in the chapter, it holds on validation and test
alike, and it has a methodological consequence that reaches beyond this project:
a comparison between state representations is meaningless for a method that does
not read its state, and reward alone cannot reveal that, because in a rising
market the state-blind policy scores well. Reporting the conditioning diagnostic
alongside returns is what made this visible, and it was added only after
identical wealth changes across a twelvefold difference in state dimension
prompted the question.

\paragraph{3. Lower drawdown by lower exposure is not risk management.} Deep Q
attains $41\%$ of buy-and-hold's return at $40\%$ of its drawdown, and the
risk-aware reward moves both figures down together. Because the ratio barely
changes and the Sharpe ratio falls, the honest description is de-risking rather
than skill. Reporting mean exposure beside every return is what makes this
checkable; without it, a table showing a threefold drawdown reduction would look
like an achievement.

\paragraph{4. Baselines and action models decide what a result means.} Three
findings in this chapter exist only because the benchmarks were constructed
carefully. The corrected buy-and-hold is $21.6\%$ stronger than the
dollar-cost-average that had been carrying its name
(Section~\ref{sec:defects}). The exposure ceiling shows that under the original
trade slice no agent could have beaten it regardless of skill
(Section~\ref{sec:res_calibration}). And reporting the slice schedule alongside
buy-and-hold is what reveals REINFORCE's near-benchmark return to be that
schedule rather than a near-miss at market timing. Weak or mislabelled baselines
would have turned each of these into a positive result.

\paragraph{5. Variance is a finding, not a nuisance.} The learned outcomes are
bimodal across seeds rather than scattered around a mean, which is why the seed
count was raised to six and why ranges are reported rather than standard errors.
It is also why the ladder's stopping rule terminates immediately: the rule
compares differences against seed spread, and the spread is larger than any
difference. A study of this design reporting a single seed could have shown
almost any conclusion it wished.

\paragraph{Scope.} A single asset, daily bars in monthly episodes, three
algorithms, a long-only three-action space, and 24 held-out months in a period
that rose overall. The negative results above are conditional on that setting, and
in particular on a reward that pays only for expected wealth. Chapter~6 sets out
which of these limits are incidental and which are load-bearing.

\paragraph{Chapter summary (plain language).} The short version is that none of
the three learners beat simply buying at the start of the month and holding, and
that the reasons are more interesting than the verdict.

Two of the three never used the market information they were given at all. They
settled on a blanket rule --- buy whenever there is cash, or never buy --- and a
diagnostic added for the purpose confirmed they were making the same choice
regardless of what the market was doing. One of them scores respectably, but only
because buying whenever possible is a decent strategy in a market that rose; it
is not evidence of skill, and the same agent ignores the cost of trading even
when that cost is made ten times larger.

The third learner did read the market, and used it to stay mostly out. It ended
up with roughly 40\% of the market's return and roughly 40\% of the market's
worst losses, which is what anyone gets by investing 40\% of their money. Asking
it explicitly to avoid losses made it invest even less, which is the cheapest way
to comply rather than the intended lesson.

The most useful discovery was not about learning at all. The original rules let
the agent commit only a tenth of its money per day, so in a twenty-one-day month
it could never become as invested as the benchmark, and could never have won
whatever it learned. Fixing that mattered more than any change to what the agent
could see.
